\NativeSlot{45}{45}{%
\ExamGroup{二、选择（将正确答案的序号填在括号里）。{\normalfont\kaishu\fontsize{10.8bp}{14bp}\selectfont （每题2分，共14分）}}

% p038_s01

\ExamQuestion{1.}{5}%
两个完全相同的三角形一定能拼成\mbox{（\hspace{1.8em}）}。\par
\ExamOptions{A. 梯形}{B. 平行四边形}{C. 长方形}{D. 正方形}
}
\NativeSlot{122.0}{44}{%
% p038_s02

\ExamQuestion{2.}{5}%
广场花坛的形状是平行四边形，相邻两条边长分别是7米和6米，这个花坛的周长是多少米？\mbox{（\hspace{1.8em}）}\par
\ExamOptions{A. 13米}{B. 14米}{C. 12米}{D. 26米}
}
\NativeSlot{199}{36}{%
% p038_s03

\ExamQuestion{3.}{5}%
一个等腰直角三角形，其中一个底角是\mbox{（\hspace{1.8em}）}。\par
\ExamOptions{A. \(90^\circ\)}{B. \(180^\circ\)}{C. \(45^\circ\)}{D. \(60^\circ\)}
}
\NativeSlot{249}{58}{%
% p038_s04

\ExamQuestion{4.}{5}%
校门口的电动伸缩门，能自由拉伸开合，它是利用了平行四边形\mbox{（\hspace{1.8em}）}的特性设计的。\par
\ExamOptions{A. 稳定性}{B. 易变形}{C. 对称性}{D. 牢固性}
\BookFigure{figures/p038_s04-01.tex}
}
\NativeSlot{350}{42}{%
% p038_s05

\ExamQuestion{5.}{5}%
下面各组数据，不能组成三角形的是\mbox{（\hspace{1.8em}）}。\par
\ExamOptions{A. 6厘米、6厘米、6厘米}{B. 3分米、3分米、6分米}{C. 9米、7米、8米}{D. 12厘米、6厘米、10厘米}
}
\NativeSlot{424}{59}{%
% p038_s06

\ExamQuestion{6.}{5}%
裁剪窗花时，把一张平行四边形红纸任意分成两个梯形，这两个梯形的\mbox{（\hspace{1.8em}）}一定相等。\par
\ExamOptions{A. 面积}{B. 周长}{C. 高}{D. 上底与下底的和}
}
\NativeSlot{524}{36}{%
% p038_s07

\ExamQuestion{7.}{5}%
一个三角形中，最小的角是\(46^\circ\)，这个三角形是\mbox{（\hspace{1.8em}）}。\par
\ExamOptions{A. 锐角三角形}{B. 直角三角形}{C. 钝角三角形}{D. 无法确定}
}
\NativeSlot{574}{36}{%
\ExamGroup{三、判断。{\normalfont\kaishu\fontsize{10.8bp}{14bp}\selectfont （每题2分，共14分）}}

% p038_s08

\ExamQuestion{1.}{5}%
判断：长方形是特殊的平行四边形，正方形是特殊的长方形。\mbox{（\hspace{1.8em}）}\par
}
\NativeSlot{623.0}{36}{%
% p038_s09

\ExamQuestion{2.}{5}%
判断：梯形的高一定比它的腰短。\mbox{（\hspace{1.8em}）}\par
}
\NativeSlot{647.5}{36}{%
% p038_s10

\ExamQuestion{3.}{5}%
判断：从平行四边形的一个顶点可以向对边作一条高，所以平行四边形有4条高。\mbox{（\hspace{1.8em}）}\par
}
\NativeSlot{698.0}{36}{%
% p038_s11

\ExamQuestion{4.}{5}%
判断：两个完全一样的梯形可以拼成一个平行四边形。\mbox{（\hspace{1.8em}）}\par
}
\NativeSlot{723.0}{36}{%
% p038_s12

\ExamQuestion{5.}{5}%
判断：剪下三角形的一个\(45^\circ\)角，剩下图形的内角和是\(135^\circ\)。\mbox{（\hspace{1.8em}）}\par
}
\NativeSlot{747.5}{36}{%
% p038_s13

\ExamQuestion{6.}{5}%
判断：三角形越大，内角和就越大。\mbox{（\hspace{1.8em}）}\par
}
\NativeSlot{773.0}{36}{%
% p038_s14

\ExamQuestion{7.}{5}%
判断：有一组对边平行的图形一定是梯形。\mbox{（\hspace{1.8em}）}\par
}
